% Response letter — LaTeX conversion target for rr-response (default draft format is Markdown).
% Derived from the RFS ETFcomove referee letters (Weikai Li's attested format).
% Ritual slots carry "% SRC:" pointers — copy the phrasing VERBATIM from references/style.md,
% adapting only the bracketed substance. Never invent phrasings.
\documentclass[12pt]{article}

\usepackage[T1]{fontenc}
\usepackage[utf8]{inputenc}
\usepackage{lmodern}
\usepackage[letterpaper,margin=1in]{geometry}
\usepackage{amsmath,amssymb,setspace,microtype}
\usepackage{booktabs,array,longtable,pdflscape}
\usepackage[authoryear,round]{natbib}
\usepackage[hidelinks]{hyperref}
\setlength{\parindent}{0pt}
\setlength{\parskip}{0.5\baselineskip}
\setlength{\emergencystretch}{2em}
\newcommand{\nothing}[1]{{}}   % on/off wrapper for appendix material

\begin{document}

\setstretch{1.5}

% ===== TITLE BLOCK (style §2 block 1) =====
\begin{center}
    \underline{\textbf{Response to Referee [N]'s Comments}}
    \par\medskip
    \textit{[Manuscript Title]}
    \par ([Journal MS \#])
    % \par [Date]   % include when journal convention expects it
\end{center}

% ===== OPENING FRAMING (style §2 block 3; §5.1 for phrasing) =====
% Beats: (a) gratitude — (b) improvement claim — (c) italics-and-response convention announced.
% Optional: global changes beyond the point-by-point (footnote, P7 pattern).
% SRC: style §5.1 — copy opener template verbatim
We are grateful for your insightful and constructive comments, [IMPROVEMENT CLAIM]. Below, we present your comments in italics, followed by our responses.\footnote{[Global changes beyond the point-by-point responses, if any — e.g., sample extension, notation change.]}

% ===== SUMMARY OF MAJOR CHANGES (style §2 block 4 — heavy revisions only) =====
% Ordered to mirror the acceptance condition (strategy §2 Pass 1, §7.1).
% Frame items as adoptions, not defenses: (a) what the referee asked + (b) what was done + (c) pointer.
\textbf{Summary of major changes in this version of the paper:}
\begin{enumerate}
    \item [Change 1 — ask, action, pointer]
    \item [Change 2 — ask, action, pointer]
\end{enumerate}

% ===== PER-COMMENT LOOP (style §2 block 6; follow the referee's OWN ordering/numbering) =====
% Preserve the referee's numbering exactly via \setcounter{enumi}{<N-1>}; keep nested structure.
% Response label: bold "Response:" (this template) OR plain prose (majority default, style §8.3) —
% pick ONE convention per letter and stay consistent.
\begin{enumerate}
    \setcounter{enumi}{0}   % = comment 1
    \item \textit{[Referee comment VERBATIM, copied from comments_extracted.md — never re-extracted]}
\end{enumerate}

% SRC: style §5.1/§5.2 — gratitude ritual + action bridge templates
% Four-beat reply: (1) thanks/agreement (2) why the point matters (3) action or new evidence
% (4) manuscript or R-table pointer. Numbers carry t-stats inline; magnitudes carry benchmarks.
\textbf{Response:} [Reply.]

% Pasting revised manuscript prose as proof of completion (style §7.8):
\begin{quote} \vspace{-2ex}
    ``[Revised manuscript text, verbatim]'' \vspace{-2ex}
\end{quote}

% Repeat the loop per comment. Multi-part replies: "First, … Second, … Finally, …";
% land takeaways with "Overall, …" (style §8.2).

% ===== CLOSING (style §2 block 9 — one ritualized sentence) =====
% SRC: style §9.10 — copy closing template verbatim
Once again, [closing ritual sentence]. The authors.

% ===== REFERENCES (style §2 block 7 — letter as free-standing document) =====
\begin{singlespace}
    \clearpage
    \bibliographystyle{mystyle}
    \bibliography{ref}
    \clearpage
\end{singlespace}

% ===== RESPONSE-ONLY EXHIBITS (style §7.5 — R-prefix, separate from manuscript numbering) =====
% Booktabs, no vertical rules; full caption + notes + t-stats.
\nothing{\clearpage
\begin{landscape}
\begin{center}
Table R1. [Caption]
\end{center}
\begingroup\setstretch{1}\footnotesize
\begin{tabular*}{\linewidth}{@{\extracolsep{\fill}}lccc}
\toprule
            & (1)   & (2)   & (3)   \\
\midrule
[Variable]  & [coef] & [coef] & [coef] \\
\bottomrule
\end{tabular*}
\endgroup
[Notes: significance stars, SE clustering, variable definitions.]
\end{landscape}}

\end{document}
